% Emacs, this is -*-latex-*-

\ex{Amostragem por import\^ancia para estimar probabilidades de cauda {\small \citep[baseado em][Exerc\'icio 3.5]{Robert2010}}}
\label{ex:importance-sampling-to-estimate-tail-probabilities}
Gostar\'iamos de usar amostragem por import\^ancia para computar a probabilidade
de que uma vari\'avel aleat\'oria Gaussiana padr\~ao $x$ assuma um valor maior
que $5$, i.e.
\begin{equation}
    \Pr(x>5) = \int_{5}^\infty \frac{1}{\sqrt{2\pi}} \exp\left(-\frac{x^2}{2}\right) dx
\end{equation}
Sabemos que a probabilidade \'e igual a
\begin{align}
    \Pr(x>5) & = 1-\int_{-\infty}^5 \frac{1}{\sqrt{2\pi}} \exp\left(-\frac{x^2}{2}\right) dx\\
    & = 1-\Phi(5)\\
    & \approx 2.87 \cdot 10^{-7}
\end{align}
onde $\Phi(.)$ \'e a fun\'c\~ao de distribui\'c\~ao acumulada de uma vari\'avel
aleat\'oria normal padr\~ao.

\begin{exenumerate}
    \item Com a fun\'c\~ao indicadora $\mathbbm{1}_{x>5}(x)$, que
    \'e igual a um se $x$ for maior que $5$ e zero caso contr\'ario, podemos
    escrever $\Pr(x>5)$ na forma da esperan\'ca
    \begin{align}
        \Pr(x>5) &= \E[ \mathbbm{1}_{x>5}(x) ],
    \end{align}
    onde a esperan\'ca \'e tomada em rela\'c\~ao \`a densidade
    $\mathcal{N}(x; 0, 1)$ de uma vari\'avel aleat\'oria normal padr\~ao,
    \begin{equation}
        \mathcal{N}(x; 0, 1) = \frac{1}{\sqrt{2\pi}} \exp\left(-\frac{x^2}{2}\right).
    \end{equation}
    Isso sugere que podemos aproximar $\Pr(x>5)$ por uma m\'edia de Monte Carlo
    \begin{equation}
        \Pr(x>5) \approx \frac{1}{n} \sum_{i=1}^n \mathbbm{1}_{x>5}(x_i), \quad \quad \quad x_i \sim  \mathcal{N}(x; 0, 1).
    \end{equation}
    Explique por que esta abordagem n\~ao funciona bem.
    
    \begin{solution}
        Nesta abordagem, n\'os essencialmente contamos quantas vezes os $x_i$
        s\~ao maiores que 5. No entanto, sabemos que a chance de que $x_i>5$
        \'e de apenas $2.87 \cdot 10^{-7}$. Ou seja, s\'o obtemos cerca de um valor acima de 5
        a cada 20 milh\~oes de simula\'c\~oes! A abordagem \'e, portanto, muito ineficiente em termos de amostras.
    \end{solution}
    
    \item Outra abordagem \'e usar amostragem por import\^ancia com uma distribui\'c\~ao de import\^ancia $q(x)$ que seja zero para $x<5$. Podemos ent\~ao escrever $\Pr(x>5)$ como
    \begin{align}
        \Pr(x>5) & = \int_{5}^\infty \frac{1}{\sqrt{2\pi}} \exp\left(-\frac{x^2}{2}\right) dx\\
        & = \int_{5}^\infty \frac{1}{\sqrt{2\pi}} \exp\left(-\frac{x^2}{2}\right) \frac{q(x)}{q(x)} dx \\
        & = \E_{q(x)} \left[\frac{1}{\sqrt{2\pi}} \exp\left(-\frac{x^2}{2}\right) \frac{1}{q(x)}\right]
    \end{align}
    e estimar $\Pr(x>5)$ como uma m\'edia amostral.
    
    Aqui usamos uma distribui\'c\~ao exponencial deslocada em $5$ para a direita. Ela possui a fdp
    \begin{equation}
        q(x) = \begin{cases}
            \exp(- (x-5)) & \text{se } x\ge 5\\
            0 & \text{caso contr\'ario}
        \end{cases}
    \end{equation}
    Para fundamentos sobre a distribui\'c\~ao exponencial, veja
    e.g.\ \url{https://en.wikipedia.org/wiki/Exponential_distribution}.
    
    Forne\'ca uma f\'ormula que aproxime $\Pr(x>5)$ como uma m\'edia amostral
    sobre $n$ amostras $x_i \sim q(x)$.
    
    \begin{solution}
        A equa\'c\~ao fornecida
        \begin{equation}
            \Pr(x>5)= \E_{q(x)} \left[\frac{1}{\sqrt{2\pi}} \exp\left(-\frac{x^2}{2}\right) \frac{1}{q(x)}\right]
        \end{equation}
        pode ser aproximada como uma m\'edia amostral da seguinte forma:
        \begin{align}
            \Pr(x>5) &\approx \frac{1}{n} \sum_{i=1}^n \frac{1}{\sqrt{2\pi}} \exp\left(-\frac{x_i^2}{2}\right) \frac{1}{q(x_i)}\\
            & = \frac{1}{n} \sum_{i=1}^n  \frac{1}{\sqrt{2\pi}} \exp\left(-\frac{x_i^2}{2} + x_i-5\right)
        \end{align}
        com $x_i \sim q(x)$.
    \end{solution}
    
    \item Compute numericamente a estimativa de import\^ancia para v\'arios tamanhos de amostra $n \in [0, 1000]$. Plote a estimativa contra o tamanho da amostra e compare com o valor real (ground truth).
    
    \begin{solution}
        A figura seguinte mostra a estimativa de amostragem por import\^ancia como uma
        fun\'c\~ao do tamanho da amostra (os n\'umeros dependem da semente aleat\'oria
        utilizada). Podemos ver que conseguimos obter uma boa estimativa j\'a com algumas
        centenas de amostras.
        \begin{center}
            \includegraphics[width=0.75 \textwidth]{IS-tail-proba}
        \end{center}
        
        O c\'odigo em Python \'e o seguinte:
        \begin{lstlisting}
            import numpy as np
            from numpy.random import default_rng
            import matplotlib.pyplot as plt
            from scipy.stats import norm
            
            n = 1000
            alpha = 5
            
            # computa a probabilidade de cauda
            p = 1-norm.cdf(alpha)
            
            # amostra da distribuicao de importancia
            rng = default_rng()
            vals = rng.exponential(scale=1, size=n) + alpha
            
            # computa a media
            def w(x):
            return 1/np.sqrt(2*np.pi)*np.exp(-x**2/2+x-alpha)
            
            Ihat = np.cumsum(w(vals))/ np.arange(1, n+1)
            
            # plot
            plt.plot(Ihat)
            plt.axhline(y=p, color="r")
            plt.xlabel("numero de amostras")
        \end{lstlisting}
        
        E o c\'odigo em Julia \'e:
        \begin{lstlisting}
            using Distributions
            using Plots
            using Statistics
            
            # computa a probabilidade de cauda
            phi(x) = cdf(Normal(0,1),x)
            alpha = 5
            p = (1-phi(alpha))
            
            # amostra da distribuicao de importancia
            n = 1000
            exprv = Exponential(1)
            x = rand(exprv, n).+alpha;
            
            # computa a aproximacao
            w(x) = 1/sqrt(2*pi)*exp(-x^2/2+x-alpha)
            
            Ihat = zeros(length(x));
            for k in 1:length(x)
            Ihat[k] = mean(w.(x[1:k]));
            end
            
            # plot
            plt=plot(Ihat, label="aproximacao");
            hline!([p], color=:red, label="valor real")
            xlabel!("numero de amostras")
        \end{lstlisting}
        
    \end{solution}
    
\end{exenumerate}


\ex{Integra\'c\~ao de Monte Carlo e amostragem por import\^ancia}
\label{ex:monte-carlo-integration-and-importance-sampling}
Uma distribui\'c\~ao Cauchy padr\~ao possui a fun\'c\~ao densidade (fdp)
\begin{equation}
    \label{eq:cauchy-pdf-alt}
    p(x) = \frac{1}{\pi}\frac{1}{1+x^2}
\end{equation}
com $x \in \mathbb{R}$. Um amigo gostaria de verificar que $\int p(x) dx = 1$, mas n\~ao sabe exatamente como resolver a integral analiticamente. Ele ent\~ao utiliza amostragem por import\^ancia e aproxima a integral como
\begin{equation}
    \int p(x) dx \approx \frac{1}{n} \sum_{i=1}^n \frac{p(x_i)}{q(x_i)} \quad \quad x_i \sim q
\end{equation}
onde $q$ \'e a densidade da distribui\'c\~ao auxiliar/de import\^ancia. Seu amigo escolhe uma densidade normal padr\~ao para $q$ e produz a seguinte figura:
\begin{center}
    \includegraphics[width=0.65\textwidth]{traces_cauchy}
\end{center}
A figura mostra duas execu\'c\~oes independentes. Em cada execu\'c\~ao, seu amigo
computa a aproxima\'c\~ao com diferentes tamanhos de amostra, incluindo subsequentemente mais e mais $x_i$ na aproxima\'c\~ao, de modo que, por exemplo, a
aproxima\'c\~ao com $n=2000$ compartilha as primeiras 1000 amostras com a
aproxima\'c\~ao que usa $n=1000$.

Seu amigo est\'a intrigado por as duas execu\'c\~oes darem resultados bastante diferentes
(que n\~ao s\~ao iguais a um) e tamb\'em por, dentro de cada execu\'c\~ao, a
estimativa depender muito do tamanho da amostra. Explique estas descobertas.

\begin{solution}
    Embora a estimativa $\hat{I}_n$
    \begin{equation}
        \hat{I}_n =  \frac{1}{n} \sum_{i=1}^n \frac{p(x_i)}{q(x_i)}
    \end{equation}
    seja n\~ao enviesada por constru\'c\~ao, temos que verificar se o seu segundo
    momento \'e finito. Caso contr\'ario, temos um estimador inv\'alido que
    se comporta de forma err\'atica na pr\'atica. A raz\~ao $w(x)$ entre $p(x)$ e
    $q(x)$ \'e igual a
    \begin{align}
        w(x) & = \frac{p(x)}{q(x)}\\
        & = \frac{\frac{1}{\pi}\frac{1}{1+x^2}}{ \frac{1}{\sqrt{2\pi}}\exp(-x^2/2)}
    \end{align}
    que pode ser simplificada para
    \begin{equation}
        w(x)  = \frac{\sqrt{2\pi}\exp(x^2/2)}{\pi(1+x^2)}.
    \end{equation}
    O segundo momento de $w(x)$ sob $q(x)$, portanto, \'e
    \begin{align}
        \E_{q(x)}\left[w(x)^2\right] & = \int_{-\infty}^\infty \frac{2\pi}{\pi^2} \frac{\exp(x^2)}{(1+x^2)^2} q(x) dx\\
        & = \int_{-\infty}^\infty \frac{2\pi}{\pi^2} \frac{\exp(x^2)}{(1+x^2)^2} \frac{1}{\sqrt{2\pi}} \exp(-x^2/2) dx\\
        & \propto \int_{-\infty}^\infty \frac{\exp(x^2/2)}{(1+x^2)^2} dx
    \end{align}
    A fun\'c\~ao exponencial cresce mais rapidamente do que qualquer polin\^omio, de
    modo que a integral torna-se arbitrariamente grande. Portanto, o segundo
    momento (e a vari\^ancia) de $\hat{I}_n$ \'e ilimitado, o que
    explica o comportamento err\'atico das curvas no gr\'afico.
    
    Uma maneira menos formal, mas mais r\'apida, de ver que, para este problema, uma
    normal padr\~ao \'e uma m\'a escolha de distribui\'c\~ao de import\^ancia \'e
    notar que sua densidade decai mais rapidamente do que a fdp de Cauchy em
    \eqref{eq:cauchy-pdf-alt}, o que significa que a fdp normal padr\~ao \'e
    ``pequena'' quando a fdp de Cauchy ainda \'e ``grande'' (veja a Figura
    \ref{fig:gauss-cauchy-logpdf-alt}). Isso leva a uma grande vari\^ancia da
    estimativa. A conclus\~ao geral \'e que a integral $\int p(x) dx$
    n\~ao deve ser aproximada com amostragem por import\^ancia com uma distribui\'c\~ao
    de import\^ancia Gaussiana.
    
    \begin{figure}[h!]
        \centering
        \includegraphics[width=0.6\textwidth]{gauss-cauchy-logpdf}
        \caption{\label{fig:gauss-cauchy-logpdf-alt} Exerc\'icio
            \ref{ex:monte-carlo-integration-and-importance-sampling}. Compara\'c\~ao
            da log-fdp de uma normal padr\~ao (azul) e da vari\'avel aleat\'oria de Cauchy
            (vermelho) para entradas positivas. A fdp de Cauchy possui caudas muito
            mais pesadas do que uma Gaussiana, de modo que a fdp Gaussiana j\'a \'e
            ``pequena'' quando a fdp de Cauchy ainda \'e ``grande''.}
    \end{figure}
\end{solution}



\ex{Amostragem por transforma\'c\~ao inversa}
\label{ex:ex1-sampling}
A fun\'c\~ao de distribui\'c\~ao acumulada (fda) $F_x(\alpha)$ de uma vari\'avel
aleat\'oria (cont\'inua ou discreta) $x$ indica a probabilidade de que $x$ assuma
valores menores ou iguais a $\alpha$,
\begin{equation}
    F_x(\alpha) = \Pr( x \le \alpha).
\end{equation}
Para vari\'aveis aleat\'orias cont\'inuas, a fda \'e definida via a integral
\begin{equation}
    F_x(\alpha) = \int_{-\infty}^{\alpha} p_x(u) \ud u,
\end{equation}
onde $p_x$ denota a fdp da vari\'avel aleat\'oria $x$ ($u$ \'e aqui uma
vari\'avel muda). Note que $F_x$ mapeia o dom\'inio de $x$ para o intervalo $[0,1]$. Por simplicidade, assumimos aqui que $F_x$ seja invers\'ivel.

Para uma vari\'avel aleat\'oria cont\'inuas $x$ com fda $F_x$, mostre que a vari\'avel aleat\'oria $y = F_x(x)$ possui distribui\'c\~ao uniforme em $[0,1]$. 

Importante: isso implica que para uma vari\'avel aleat\'oria $y$ que \'e
uniformemente distribu\'ida em $[0,1]$, a vari\'avel aleat\'oria transformada
$F_x^{-1}(y)$ possui fda $F_x$. Isso d\'a origem a um m\'etodo chamado
``amostragem por transforma\'c\~ao inversa'' para gerar $n$ amostras iid
de uma vari\'avel aleat\'oria $x$ com fda $F_x$.
%
Dada uma fda alvo $F_x$, o m\'etodo consiste em:
\begin{itemize}
    \item calcular a inversa $F_x^{-1}$
    \item amostrar $n$ vari\'aveis aleat\'orias iid uniformemente distribu\'idas em $[0,1]$: $y^{(i)} \sim \mathcal{U}(0,1)$, $i=1, \ldots, n$.
    \item transformar cada amostra por $F_x^{-1}$: $x^{(i)} = F_x^{-1}(y^{(i)})$, $i=1, \ldots, n$.
\end{itemize}
Pela constru\'c\~ao do m\'etodo, os $x^{(i)}$ s\~ao $n$ amostras iid de $x$.

\begin{solution}
    Come\'camos com a fun\'c\~ao de distribui\'c\~ao acumulada (fda) $F_y$ para $y$,
    \begin{align}
        F_y(\beta) & = \Pr( y \le \beta).
    \end{align}
    Como $F_x(x)$ mapeia $x$ para $[0,1]$, $F_y(\beta)$ \'e zero para $\beta < 0$ e um para $\beta > 1$. Em seguida, consideramos $\beta \in [0,1]$.
    
    Seja $\alpha$ o valor de $x$ que $F_x$ mapeia para $\beta$, i.e.\ $F_x(\alpha) = \beta$, o que significa $\alpha = F_x^{-1}(\beta)$. Como $F_x$ \'e uma fun\'c\~ao n\~ao decrescente, temos
    \begin{equation}
        F_y(\beta) = \Pr( y \le \beta) = \Pr( F_x(x) \le \beta) = \Pr(x \le F_x^{-1}(\beta)) = \Pr( x \le \alpha) = F_x(\alpha).
    \end{equation}
    Como $\alpha = F_x^{-1}(\beta)$, obtemos
    \begin{equation}
        F_y(\beta)  =  F_x( F_x^{-1}(\beta) ) = \beta
    \end{equation}
    A fda $F_y$ \'e, portanto, dada por
    \begin{equation}
        F_y(\beta) = \begin{cases}
            0 & \text{se } \beta < 0\\
            \beta & \text{se } \beta \in [0,1]\\
            1 & \text{se } \beta > 1
        \end{cases}
    \end{equation}
    que \'e a fda de uma vari\'avel aleat\'oria uniforme em $[0,1]$. Logo, $y=F_x(x)$ \'e uniformemente distribu\'ida em $[0,1]$.
    
\end{solution}



\ex{Amostragem da distribui\'c\~ao exponencial}
A distribui\'c\~ao exponencial possui a densidade
\begin{equation}
    p(x; \lambda) = \begin{cases} \lambda \exp(-\lambda x) & x \ge 0\\ 0
        & x<0,
    \end{cases}
\end{equation}
onde $\lambda$ \'e um par\^ametro da distribui\'c\~ao. Use amostragem por transforma\'c\~ao inversa para gerar $n$ amostras iid de $p(x; \lambda)$.

\begin{solution}
    Primeiro calculamos a fun\'c\~ao de distribui\'c\~ao acumulada.
    \begin{align}
        F_x(\alpha) & = \Pr(x\le \alpha)\\
        & = \int_0^\alpha  \lambda \exp(-\lambda x)\\
        & = -\exp(-\lambda x)\big|_0^\alpha\\
        & = 1-\exp(-\lambda \alpha)
    \end{align}
    Sua inversa \'e obtida resolvendo
    \begin{equation}
        y = 1-\exp(-\lambda x)
    \end{equation}
    para $x$, o que resulta em:
    \begin{align}
        \exp(-\lambda x) &= 1-y\\
        -\lambda x &= \log(1-y)\\
        x &= \frac{-\log(1-y)}{\lambda}
    \end{align}
    Para gerar amostras $x^{(i)}  \sim p(x; \lambda)$, primeiro amostramos $y^{(i)} \sim U(0,1)$ e ent\~ao definimos
    \begin{equation}
        x^{(i)} = \frac{-\log(1-y^{(i)})}{\lambda}.
    \end{equation}
    
    A amostragem por transforma\'c\~ao inversa pode ser usada para gerar amostras
    de muitas distribui\'c\~oes padrão. Por exemplo, ela permite gerar
    vari\'aveis aleat\'orias Gaussianas a partir de vari\'aveis aleat\'orias uniformemente
    distribu\'idas. O m\'etodo \'e chamado de transforma\'c\~ao de Box-Muller, veja
    e.g. \url{https://en.wikipedia.org/wiki/Box-Muller_transform}. Como
    gerar as amostras requeridas da distribui\'c\~ao uniforme \'e um campo de
    pesquisa por si s\'o, veja
    e.g. \url{https://en.wikipedia.org/wiki/Random_number_generation}
    e \citep[Cap\'itulo 3]{Owen2013}.
    
\end{solution}


\ex{Amostragem de uma distribui\'c\~ao de Laplace}
\label{ex:sampling-from-a-Laplace-distribution}
Uma vari\'avel aleat\'oria de Laplace $x$ de m\'edia zero e vari\^ancia um possui a densidade $p(x)$
\begin{equation}
    p(x) = \frac{1}{\sqrt{2}} \exp\left(-\sqrt{2}|x|\right) \quad \quad x \in \mathbb{R}.  
\end{equation}
Use amostragem por transforma\'c\~ao inversa para gerar $n$ amostras iid de $x$.

\begin{solution}
    A tarefa principal \'e computar a fun\'c\~ao de distribui\'c\~ao acumulada (fda) $F_x$ de $x$ e sua inversa. A fda \'e, por defini\'c\~ao,
    \begin{align}
        F_x(\alpha) & = \int_{-\infty}^\alpha  \frac{1}{\sqrt{2}} \exp\left(-\sqrt{2}|u|\right) \ud u.
    \end{align}
    Consideramos primeiro o caso em que $\alpha \le 0$. Como $-|u| = u$ para $u\le 0$, temos 
    \begin{align}
        F_x(\alpha) & = \int_{-\infty}^\alpha  \frac{1}{\sqrt{2}} \exp\left(\sqrt{2}u\right) \ud u   \\
        & = \frac{1}{2} \exp\left(\sqrt{2}u\right) \bigg |_{-\infty}^\alpha\\
        & =  \frac{1}{2} \exp\left(\sqrt{2} \alpha\right).
    \end{align}
    Para $\alpha > 0$, temos
    \begin{align}
        F_x(\alpha) & = \int_{-\infty}^\alpha  \frac{1}{\sqrt{2}} \exp\left(-\sqrt{2}|u|\right) \ud u   \\
        &= 1 - \int_{\alpha}^\infty  \frac{1}{\sqrt{2}} \exp\left(-\sqrt{2}|u|\right) \ud u        
    \end{align}
    onde usamos o fato de que a fdp deve integrar um. Para valores de $u>0$, $-|u| = -u$, de modo que
    \begin{align}
        F_x(\alpha) & = 1 - \int_{\alpha}^\infty  \frac{1}{\sqrt{2}} \exp\left(-\sqrt{2} u \right) \ud u  \\
        &= 1 + \frac{1}{2}  \exp\left(-\sqrt{2} u \right) \bigg |_{\alpha}^{\infty}\\
        &= 1 - \frac{1}{2}  \exp\left(-\sqrt{2} \alpha \right).
    \end{align}
    No total, para $\alpha \in \mathbb{R}$, temos ent\~ao
    \begin{equation}
        F_x(\alpha) = \begin{cases}
            \frac{1}{2} \exp\left(\sqrt{2} \alpha\right) & \text{se } \alpha \le 0\\
            1 - \frac{1}{2}  \exp\left(-\sqrt{2} \alpha\right) & \text{se } \alpha > 0
        \end{cases}
    \end{equation}
    A Figura \ref{fig:Fx} visualiza $F_x(\alpha)$.
    \begin{figure}[h!]
        \centering
        \includegraphics[width = 0.6 \textwidth]{cdf-laplace}
        \caption{\label{fig:Fx} A fun\'c\~ao de distribui\'c\~ao acumulada $F_x(\alpha)$ para uma vari\'avel aleat\'oria com distribui\'c\~ao de Laplace.}
    \end{figure}
    
    Como a figura sugere, existe uma inversa \'unica para $y = F_x(\alpha)$. Para $y \le 1/2$, temos
    \begin{align}
        y & =  \frac{1}{2} \exp\left(\sqrt{2} \alpha\right)\\
        \log(2y) & = \sqrt{2} \alpha\\
        \alpha & = \frac{1}{\sqrt{2}} \log (2y)
    \end{align}
    Para $y > 1/2$, temos
    \begin{align}
        y & =  1 - \frac{1}{2}  \exp\left(-\sqrt{2} \alpha \right)\\
        -y & =  -1 + \frac{1}{2}  \exp\left(-\sqrt{2} \alpha \right)\\
        1-y &= \frac{1}{2}  \exp\left(-\sqrt{2} \alpha \right) \\
        \log( 2-2y) & = -\sqrt{2} \alpha \\
        \alpha & = -\frac{1}{\sqrt{2}} \log(2-2y)
    \end{align}
    A fun\'c\~ao $y \mapsto g(y)$ que ocorre no logaritmo em ambos os casos \'e
    \begin{equation}
        g(y) = \begin{cases}
            2 y & \text{se } y\le \frac{1}{2}\\
            2-2y & \text{se } y > \frac{1}{2}
        \end{cases}.
    \end{equation}
    Ela \'e mostrada abaixo e pode ser escrita de forma mais compacta como $g(y) = 1-2|y-1/2|$.
    \begin{figure}[h]
        \centering
        \includegraphics[width = 0.6 \textwidth]{g}
    \end{figure}
    
    Podemos assim escrever a inversa $F_x^{-1}(y)$ da fda $y=F_x(\alpha)$ como
    \begin{equation}
        F_x^{-1}(y) = -\text{sign}\left(y-\frac{1}{2}\right) \frac{1}{\sqrt{2}} \log\left[  1-2\big|y-\frac{1}{2}\big| \right].
    \end{equation}
    Para gerar $n$ amostras iid de $x$, primeiro geramos $n$
    amostras iid $y^{(i)}$ que s\~ao uniformemente distribu\'idas em $[0,1]$,
    e ent\~ao computamos para cada uma $F_x^{-1}(y^{(i)})$. As propriedades da
    amostragem por transforma\'c\~ao inversa garantem que os $x^{(i)}$,
    \begin{equation}
        x^{(i)} = F_x^{-1}(y^{(i)}),
    \end{equation}
    sejam independentes e possuam distribui\'c\~ao de Laplace.
\end{solution}

\ex{Amostragem por rejei\'c\~ao {\small \citep[baseado em][Exerc\'icio 2.8]{Robert2010}}}
\label{ex:rejection-sampling-alt}
A maioria dos ambientes de computa\'c\~ao fornece fun\'c\~oes para amostrar de uma distribui\'c\~ao
normal padr\~ao. Algoritmos populares incluem a transforma\'c\~ao de
Box-Muller, veja
e.g. \url{https://en.wikipedia.org/wiki/Box-Muller_transform}. Aqui
usamos amostragem por rejei\'c\~ao para amostrar de uma distribui\'c\~ao normal
padr\~ao com densidade $p(x)$ usando uma distribui\'c\~ao de Laplace como nossa
distribui\'c\~ao de proposta/auxiliar.

A densidade $q(x)$ de uma distribui\'c\~ao de Laplace de m\'edia zero com vari\^ancia
$2b^2$ \'e
\begin{equation}
    q(x;b) = \frac{1}{2b}\exp\left(-\frac{|x|}{b}\right).
\end{equation}
Podemos amostrar dela amostrando uma vari\'avel de Laplace com vari\^ancia 1
como no Exerc\'icio \ref{ex:sampling-from-a-Laplace-distribution} e ent\~ao
escalando a amostra por $\sqrt{2}b$.

A amostragem por rejei\'c\~ao ent\~ao repete os seguintes passos:
\begin{itemize}
    \item Gerar $x \sim q(x; b)$
    \item Aceitar $x$ com probabilidade $f(x) = \frac{1}{M}
    \frac{p(x)}{q(x)}$, i.e. gerar $u \sim U(0,1)$ e aceitar $x$
    se $u \le f(x)$.
\end{itemize}

\begin{exenumerate}
    \item Compute a raz\~ao $M(b) = \max_x \frac{p(x)}{q(x; b)}$.
    \begin{solution}
        Pelas defini\'c\~oes da fdp $p(x)$ de uma normal padr\~ao e da fdp $q(x; b)$ da distribui\'c\~ao de Laplace, temos
        \begin{align}
            \frac{p(x)}{q(x;b)} &= \frac{ \frac{1}{\sqrt{2\pi}} \exp(-x^2/2)}{\frac{1}{2b}\exp(-|x|/b)}\\
            & = \frac{2b}{\sqrt{2\pi}} \exp(-x^2/2 + |x|/b)
        \end{align}
        A raz\~ao \'e sim\'etrica em $x$. Al\'em disso, como a fun\'c\~ao
        exponencial \'e estritamente crescente, podemos encontrar o maximizador de
        $-x^2/2+x/b$ para $x\ge 0$ para determinar o maximizador de
        $M(b)$. Com $g(x) = -x^2/2+x/b$, temos
        \begin{align}
            g'(x) &= -x + 1/b\\
            g''(x) &= -1
        \end{align}
        O ponto cr\'itico (para o qual a primeira derivada \'e zero) \'e $x = 1/b$ e, como a segunda derivada \'e negativa para todo $x$, o ponto \'e um m\'aximo. A raz\~ao m\'axima $M(b)$, portanto, \'e
        \begin{align}
            M(b)  & = \frac{2b}{\sqrt{2\pi}} \exp\left(-x^2/2 + |x|/b\right)\Big|_{x=1/b}\\
            & = \frac{2b}{\sqrt{2\pi}} \exp\left(-1/(2b^2) + 1/b^2\right)\\
            & = \frac{2b}{\sqrt{2\pi}} \exp\left(1/(2b^2)\right)
        \end{align}
    \end{solution}
    
    \item Como voc\^e deve escolher $b$ para maximizar a probabilidade de aceita\'c\~ao?
    
    \begin{solution}
        A probabilidade de aceita\'c\~ao \'e $1/M$. Logo, para maximiz\'a-la,
        temos que escolher $b$ tal que $M(b)$ seja m\'inimo. Computamos as derivadas
        \begin{align}
            M'(b) & = \frac{2}{\sqrt{2\pi}} \exp(1/(2b^2)) - \frac{2b}{\sqrt{2\pi}} \exp(1/(2b^2)) b^{-3}\\
            & = \frac{2}{\sqrt{2\pi}} \exp(1/(2b^2)) - \frac{2}{\sqrt{2\pi}} \exp(1/(2b^2)) b^{-2}\\
            & = \frac{2}{\sqrt{2\pi}} \exp(1/(2b^2))(1-b^{-2})\\
            M''(b) & = -b^{-3}\frac{2}{\sqrt{2\pi}} \exp(1/(2b^2))(1-b^{-2})+2b^{-3}\frac{2}{\sqrt{2\pi}} \exp(1/(2b^2))\\
        \end{align}
        Igualando a primeira derivada a zero, obtemos
        \begin{align}
            \frac{2}{\sqrt{2\pi}} \exp(1/(2b^2)) &= \frac{2}{\sqrt{2\pi}} \exp(1/(2b^2)) b^{-2}\\
            1 & = b^{-2}
        \end{align}
        Portanto, o $b$ \'otimo \'e $b=1$. A segunda derivada em $b=1$ \'e
        \begin{align}
            M''(1) & = 2 \frac{2}{\sqrt{2\pi}} \exp(1/2)
        \end{align}
        que \'e positiva, de modo que $b=1$ \'e um m\'inimo. O menor valor de $M$ \'e, portanto,
        \begin{align}
            M(1) & =  \frac{2b}{\sqrt{2\pi}} \exp(1/(2b^2))\Big |_{b=1} \\
            & =  \frac{2}{\sqrt{2\pi}} \exp(1/2)\\
            & =  \sqrt{\frac{2 e}{\pi}}
        \end{align}
        onde $e = \exp(1)$. A probabilidade de aceita\'c\~ao m\'axima, portanto, \'e
        \begin{align}
            \frac{1}{\min_b M(b)} & = \sqrt{\frac{\pi}{2e}}\\
            & \approx 0.76
        \end{align}
        Isso significa que, para cada amostra $x$ gerada de $q(x; 1)$, h\'a
        uma chance de $0.76$ de que ela seja aceita. Em outras palavras, para cada
        amostra aceita, precisamos gerar $1/0.76 = 1.32$ amostras de
        $q(x; 1)$.
        
        A vari\^ancia da distribui\'c\~ao de Laplace para $b=1$ \'e igual a 2. Assim,
        a vari\^ancia da distribui\'c\~ao auxiliar \'e maior (o dobro)
        que a vari\^ancia da distribui\'c\~ao que gostar\'iamos de amostrar. 
    \end{solution}
    \item Suponha que voc\^e amostre de $p(x_1, \ldots, x_d) = \prod_{i=1}^d
    p(x_i)$ usando $q(x_1, \ldots, x_d) = \prod_{i=1}^d q(x_i; b)$ como
    distribui\'c\~ao auxiliar sem explorar quaisquer independ\^encias. Como
    a probabilidade de aceita\'c\~ao escala como uma fun\'c\~ao de $d$? Voc\^e pode
    denotar a probabilidade de aceita\'c\~ao no caso de $d=1$ por $A$.
    
    \begin{solution}
        Temos que determinar a raz\~ao m\'axima
        \begin{equation}
            M_d = \max_{x_1, \ldots, x_d} \frac{p(x_1, \ldots, x_d)}{q(x_1, \ldots, x_d)}
        \end{equation}
        Substituindo a fatoriza\'c\~ao, obtemos
        \begin{align}
            M_d &= \max_{x_1, \ldots, x_d} \prod_{i=1}^d \frac{p(x_i)}{q(x_i)}\\
            & = \prod_{i=1}^d \underbrace{\max_{x_i} \frac{p(x_i)}{q(x_i)}}_{M_1 = 1/A}\\
            & = \prod_{i=1}^d \frac{1}{A}\\
            & = \frac{1}{A^d}
        \end{align}
        Portanto, a probabilidade de aceita\'c\~ao \'e
        \begin{equation}
            \frac{1}{M_d} = A^d
        \end{equation}
        Note que $A \le 1$ pois \'e uma probabilidade. Isso significa que, a menos que $A=1$,
        temos uma probabilidade de aceita\'c\~ao que decai exponencialmente com o
        n\'umero de dimens\~oes se as distribui\'c\~oes alvo e auxiliar
        se fatorizarem e n\~ao explorarmos as independ\^encias.
    \end{solution}
\end{exenumerate}


\ex{Amostragem de uma m\'aquina de Boltzmann restrita}

A m\'aquina de Boltzmann restrita (RBM) \'e um modelo para vari\'aveis bin\'arias
$\v =(v_1, \ldots, v_n)^\top$ e $\h=(h_1, \ldots, h_m)^\top$ que
afirma que a distribui\'c\~ao conjunta de $(\v,\h)$ pode ser descrita pela
fun\'c\~ao de massa de probabilidade
\begin{equation}
    p(\v,\h) \propto \exp\left( \v^\top \W \h + \a^\top\v + \b^\top\h \right),
\end{equation}
onde $\W$ \'e uma matriz $n \times m$, e $\a$ e $\b$ vetores de tamanho
$n$ e $m$, respectivamente. Tanto os $v_i$ quanto os $h_i$ assumem
valores em $\{0,1\}$. Os $v_i$ s\~ao chamados de vari\'aveis ``vis\'iveis''
j\'a que s\~ao assumidas como observadas, enquanto os $h_i$ s\~ao as vari\'aveis ocultas,
j\'a que se assume que n\~ao podemos medi-las.

Explique como usar a amostragem de Gibbs para gerar amostras da
marginal $p(\v)$,
\begin{equation}
    p(\v) = \frac{\sum_{\h}  \exp\left( \v^\top \W \h + \a^\top\v + \b^\top\h \right)}{\sum_{\h, \v}  \exp\left( \v^\top \W \h + \a^\top\v + \b^\top\h \right)},
\end{equation}
para quaisquer valores dados de $\W$, $\a$, e $\b$.

\emph{Dica:} Voc\^e pode usar que
\begin{align}
    p(\h | \v) &= \prod_{i=1}^m p(h_i | \v), &  p(h_i=1 | \v) & =  \frac{1}{ 1+  \exp\left(- \sum_j v_j W_{ji} - b_i \right)}, \\
    p(\v | \h) &= \prod_{i=1}^n p(v_i | \h), &  p(v_i = 1 | \h) &= \frac{1}{1+\exp\left(-\sum_j W_{ij} h_j-a_i\right)}.
\end{align}


\begin{solution}
    Para gerar amostras $\v^{(k)}$ de $p(\v)$, geramos
    amostras $(\v^{(k)}, \h^{(k)})$ de $p(\v, \h)$ e ent\~ao ignoramos os
    $\h^{(k)}$.
    
    A amostragem de Gibbs \'e um m\'etodo MCMC para produzir uma sequ\^encia de amostras
    $\x^{(1)}, \x^{(2)}, \x^{(3)}, \ldots$ que seguem uma fdp/fmp $p(\x)$
    (se a cadeia for executada por tempo suficiente). Assumindo que $\x$ seja
    $d$-dimensional, geramos a pr\'oxima amostra $\x^{(k+1)}$ na
    sequ\^encia a partir da amostra anterior $\x^{(k)}$ da seguinte forma:
    \begin{enumerate}
        \item escolhendo (aleatoriamente) um \'indice $i \in \{1, \ldots, d \}$
        \item amostrando $x_i^{(k+1)}$ de $p(x_i \mid \x^{(k)}_{\backslash
            i})$ onde $\x^{(k)}_{\backslash i}$ \'e o vetor $\x$ com $x_i$
        removido, i.e. $\x^{(k)}_{\backslash i}
        =({x}_1^{(k)},\ldots,{x}_{i-1}^{(k)},{x}_{i+1}^{(k)},\ldots,{x}_d^{(k)})$
        \item definindo $\x^{(k+1)} =({x}_1^{(k)},\ldots,{x}_{i-1}^{(k)},x_i^{(k+1)},{x}_{i+1}^{(k)},\ldots,{x}_d^{(k)})$.
    \end{enumerate}
    Para a RBM, a tupla $(\h,\v)$ corresponde a $\x$, de modo que um $x_i$
    nos passos acima pode ser tanto uma vari\'avel oculta $h_i$ quanto uma vis\'ivel $v_i$. Logo
    \begin{equation}
        p(x_i \mid \x_{\backslash i}) = \begin{cases}
            p(h_i \mid \h_{\backslash i}, \v) & \text{se $x_i$ for uma vari\'avel oculta $h_i$}\\
            p(v_i \mid \v_{\backslash i}, \h) & \text{se $x_i$ for uma vari\'avel vis\'ivel $v_i$}
        \end{cases}
    \end{equation}
    ($\h_{\backslash i}$ denota o vetor $\h$ com o elemento $h_i$
    removido, e equivalentemente para $\v_{\backslash i}$)
    
    Para computar as condicionais no lado direito, usamos a dica:
    \begin{align}
        p(\h | \v) &= \prod_{i=1}^m p(h_i | \v), &  p(h_i=1 | \v) & =  \frac{1}{ 1+  \exp\left(- \sum_j v_j W_{ji} - b_i \right)}, \\
        p(\v | \h) &= \prod_{i=1}^n p(v_i | \h), &  p(v_i = 1 | \h) &= \frac{1}{1+\exp\left(-\sum_j W_{ij} h_j-a_i\right)}.
    \end{align}
    Dadas as independ\^encias entre as ocultas dadas as vis\'iveis e vice-versa, temos
    \begin{align}
        p(h_i \mid \h_{\backslash i}, \v) & = p(h_i \mid \v) &   p(v_i \mid \v_{\backslash i}, \h) & = p(v_i \mid \h)
    \end{align}
    de modo que as express\~ao para $p(h_i=1 | \v)$ e $p(v_i = 1 | \h)$ nos permitem implementar o amostrador de Gibbs.
    
    Dadas as independ\^encias, faz sentido, ademais, amostrar as vari\'aveis $\h$
    e $\v$ em blocos: primeiro amostramos todas as $h_i$ dado
    $\v$, e ent\~ao todas as $v_i$ dado as $\h$ (ou vice-versa). Isso
    tamb\'em \'e conhecido como amostragem de Gibbs em blocos.
    
    Em resumo, dada uma amostra $(\h^{(k)},\v^{(k)})$, geramos ent\~ao a pr\'oxima amostra
    $(\h^{(k+1)},\v^{(k+1)})$ na sequ\^encia da seguinte forma:
    \begin{itemize}
        \item Para todo $h_i$, $i=1, \ldots, m$:
        \begin{itemize}
            \item computar $p^h_i = p(h_i=1 | \v^{(k)})$
            \item amostrar $u_i$ de uma distribui\'c\~ao uniforme em $[0,1]$ e definir $h^{(k+1)}_i$ como 1 se $u_i \le p^h_i$.
        \end{itemize}
        \item Para todo $v_i$, $i=1, \ldots, n$:
        \begin{itemize}
            \item computar $p^v_i = p(v_i=1 | \h^{(k+1)})$
            \item amostrar $u_i$ de uma distribui\'c\~ao uniforme em $[0,1]$ e definir $v^{(k+1)}_i$ como 1 se $u_i \le p^v_i$.
        \end{itemize}
    \end{itemize}
    Como passo final, ap\'os amostrar $S$ pares $(\h^{(k)},\v^{(k)})$, $k=1,
    \ldots, S$, o conjunto de vis\'iveis $\v^{(k)}$ forma amostras da
    marginal $p(\v)$.
\end{solution}



\ex{Infer\^encia b\'asica por Markov chain Monte Carlo} \label{q:basic-mcmc-alt}
Este exerc\'icio trata de amostragem e infer\^encia aproximada por Markov
chain Monte Carlo (MCMC). O MCMC pode ser usado para obter
amostras de uma distribui\'c\~ao de probabilidade, por exemplo, uma distribui\'c\~ao
posterior. As amostras representam a distribui\'c\~ao de forma aproximada, como
ilustrado na Figura~\ref{fig:mcmc_approx_alt}, e podem ser usadas para
aproximar esperan\'cas.

Denotamos a densidade de uma Gaussiana de m\'edia zero com vari\^ancia $\sigma^2$ por
$\normal(x; \mu, \sigma^2)$, i.e.\ 
\begin{equation}
    \normal(x; \mu, \sigma^2) = \frac{1}{\sqrt{2\pi \sigma^2}}\exp\left(-\frac{(x-\mu)^2}{2\sigma^2}\right)
\end{equation}

\begin{figure}[!h]
    \centering
    \subfloat[Densidade real]{
        \includegraphics[width=0.45\textwidth]{density}
        \label{sfig:density-alt}}
    \subfloat[Densidade representada por $10.000$ amostras.]{
        \includegraphics[width=0.45\textwidth]{scatter}
        \label{sfig:scatter-alt}
    }
    \caption{Densidade e amostras de $p(x,y) = \normal(x; 0,1)\normal(y;0,1)$.}
    \label{fig:mcmc_approx_alt}
\end{figure}

Considere um vetor de $d$ vari\'aveis aleat\'orias $\thetab = (\theta_1, \dots, \theta_d)$ e alguns dados observados $\data$.
%
Em muitos casos, estamos interessados em computar esperan\'cas sob
a distribui\'c\~ao posterior $p(\thetab \mid \data)$, e.g.\ 
\begin{equation}
    \E_{p(\thetab \mid \data)}\left[g(\thetab)\right] = \int g(\thetab) p(\thetab \mid \data) \mathrm{d} \thetab
\end{equation}
para alguma fun\'c\~ao $g(\thetab)$. Se $d$ for pequeno, por exemplo $d \le 3$, m\'etodos
num\'ericos determin\'isticos podem ser usados para aproximar a integral com alta precis\~ao, veja
e.g.\ \url{https://en.wikipedia.org/wiki/Numerical_integration}. Mas para
dimens\~oes maiores, esses m\'etodos geralmente n\~ao s\~ao mais aplic\'aveis. A
esperan\'ca, contudo, pode ser aproximada como uma m\'edia amostral se tivermos amostras
$\thetab^{(i)}$ de $p(\thetab \mid \data)$:
\begin{equation}
    \E_{p(\thetab \mid \data)}\left[g(\thetab)\right] \approx \frac{1}{S}\sum_{i=1}^{S} g(\thetab^{(i)})
\end{equation}
Note que em m\'etodos MCMC, as amostras $\{ \thetab^{(1)}, \ldots, \thetab^{(S)} \}$ usadas na
aproxima\'c\~ao acima tipicamente n\~ao s\~ao estatisticamente independentes.

Metropolis-Hastings \'e um algoritmo MCMC que gera amostras de uma
distribui\'c\~ao $p(\thetab)$, onde $p(\thetab)$ pode ser qualquer distribui\'c\~ao
nos par\^ametros (e n\~ao apenas posteriors). O algoritmo \'e iterativo
e, na itera\'c\~ao $t$, utiliza:
\begin{itemize}
    \item uma distribui\'c\~ao de proposta $q(\thetab; \thetab^{(t)})$, parametrizada pelo
    estado atual da cadeia de Markov, i.e.\ $\thetab^{(t)}$;
    \item uma fun\'c\~ao $p^*(\thetab)$, que \'e proporcional a $p(\thetab)$. Em outras
    palavras, $p^*(\thetab)$ n\~ao est\'a normalizada e a densidade normalizada
    $p(\thetab)$ \'e:
    \begin{equation} p(\thetab) = \frac{p^*(\thetab)}{\int p^*(\thetab) \mathrm{d}\thetab}.
    \end{equation}
\end{itemize} 
Para todas as tarefas neste exerc\'icio, trabalhamos com uma distribui\'c\~ao de proposta
Gaussiana $q(\thetab; \thetab^{(t)})$, cuja m\'edia \'e a amostra anterior na cadeia de
Markov, e cuja vari\^ancia \'e $\epsilon^2$. Ou seja, na itera\'c\~ao $t$ do nosso
algoritmo Metropolis-Hastings,
\begin{equation}
    q(\thetab; \thetab^{(t-1)}) = \prod_{k=1}^d\normal(\theta_k; \theta_k^{(t-1)}, \epsilon^2).
    \label{eq:gauss-prop-alt}
\end{equation}
Quando utilizado com esta distribui\'c\~ao de proposta, o algoritmo \'e
chamado de algoritmo Random Walk Metropolis-Hastings.
%


\begin{exenumerate}
    \item Leia a Se\'c\~ao 27.4 de \citet{Barber2012} para se familiarizar com o algoritmo Metropolis-Hastings.
    
    \item\label{qpt:mh-alt} Escreva uma fun\'c\~ao \lstinline|mh| implementando o
    algoritmo Metropolis Hastings, como dado no Algoritmo 27.3 em
    \citet{Barber2012}, usando a distribui\'c\~ao de proposta Gaussiana em \eqref{eq:gauss-prop-alt} acima. A
    fun\'c\~ao deve receber como argumentos:
    \begin{itemize}
        \item \lstinline{p_star}: uma fun\'c\~ao sobre $\thetab$ que \'e proporcional \`a densidade de interesse $p(\thetab)$;
        \item \lstinline{param_init}: a amostra inicial --- um valor para $\thetab$ de onde a cadeia de Markov come\'ca;
        \item \lstinline{num_samples}: o n\'umero $S$ de amostras a gerar;
        \item \lstinline{vari}: a vari\^ancia $\epsilon^2$ para a distribui\'c\~ao de proposta Gaussiana $q$;
    \end{itemize}
    e retornar $[\thetab^{(1)}, \dots, \thetab^{(S)}]$ --- uma lista de $S$ amostras de $p(\thetab) \propto p^*(\thetab)$. Por exemplo:
    \begin{lstlisting}
        def mh(p_star, param_init, num_samples=5000, vari=1.0):
        # seu codigo aqui
        return samples
    \end{lstlisting}
    
    
    
    \begin{solution}
        
        
        Abaixo est\'a uma implementa\'c\~ao em Python.
        \begin{lstlisting}
            def mh(p_star, param_init, num_samples=5000, vari=1.0):
            x = []
            
            x_current = param_init
            for n in range(num_samples):
            
            # proposta
            x_proposed = multivariate_normal.rvs(mean=x_current, cov=vari)
            
            # passo MH
            a = multivariate_normal.pdf(x_current, mean=x_proposed, cov=vari) * p_star(x_proposed)
            a = a / (multivariate_normal.pdf(x_proposed, mean=x_current, cov=vari) * p_star(x_current))
            
            # aceitar ou nao 
            if a >= 1:
            x_next = np.copy(x_proposed)
            elif uniform.rvs(0, 1) < a:
            x_next = np.copy(x_proposed)
            else:
            x_next = np.copy(x_current)
            
            # manter registro
            x.append(x_next)
            x_current = x_next
            
            return x
        \end{lstlisting}
        %
        Como estamos usando uma distribui\'c\~ao de proposta sim\'etrica, $q(\thetab \mid \thetab^*) = q(\thetab^* \mid \thetab)$, e poder\'iamos simplificar o algoritmo tendo $a = \frac{p^*(\thetab^*)}{p^*(\thetab)}$, onde $\thetab$ \'e a amostra atual e $\thetab^*$ \'e a amostra proposta.
        
        Na pr\'atica, \'e desej\'avel implementar a fun\'c\~ao no dom\'inio logar\'itmico para evitar problemas num\'ericos. Ou seja, em vez de $p^*$, \lstinline|mh| aceitar\'a como argumento $\log p^*$, e $a$ ser\'a calculado como: 
        $$a = (\log q(\thetab \mid \thetab^*) + \log p^*(\thetab^*)) - (\log q(\thetab^* \mid \thetab) + \log p^*(\thetab))$$
    \end{solution}
    
    \item \label{qpt:mh_test-alt} Teste seu algoritmo amostrando $5.000$
    amostras de $p(x, y) = \normal(x; 0, 1)\normal(y; 0,
    1)$. Inicialize em $(x=0, y=0)$ e use $\epsilon^2=1$.
    %
    Gere um gr\'afico de dispers\~ao das amostras obtidas. O gr\'afico deve ser
    semelhante \`a Figura~\ref{sfig:scatter-alt}. Destaque apenas as primeiras 20 amostras. 
    Somente estas 20 amostras aproximam adequadamente a densidade real?
    
    Amostre outros $5.000$ pontos de $p(x, y) = \normal(x; 0,
    1)\normal(y; 0, 1)$ usando \lstinline|mh| com $\epsilon^2=1$, mas desta
    vez inicialize em $(x=7,y=7)$. Gere um gr\'afico de dispers\~ao das amostras colhidas
    e destaque as primeiras 20 amostras. Se tudo correu conforme o esperado, seu gr\'afico
    provavelmente mostra uma ``trilha'' de amostras, come\'cando em
    $(x=7, y=7)$ e aproximando-se lentamente da regi\~ao do espa\'co onde se encontra a maior
    parte da massa de probabilidade.
    
    
    \begin{solution}
        
        \begin{figure}[!h]
            \captionsetup{type=figure}
            \centering
            \subfloat[\label{sfig:00-alt} Iniciando a cadeia em $(0,0)$.]{
                \includegraphics[width=0.45\textwidth]{00_warmup}}
            \subfloat[\label{sfig:77-alt} Iniciando a cadeia em $(7,7)$]{
                \includegraphics[width=0.45\textwidth]{77_warmup}}
            \caption{\label{fig:initial-alt} $5.000$ amostras de $\normal(x; 0, 1)\normal(y; 0, 1)$ (azul), com as primeiras $20$ amostras destacadas (vermelho). Colhidas usando Metropolis-Hastings com diferentes pontos de partida.}
        \end{figure}
        
        
        A Figura~\ref{fig:initial-alt} mostra os dois gr\'aficos de dispers\~ao de sorteios de $\normal(x; 0, 1)\normal(y; 0, 1)$:
        \begin{itemize}
            \item A Figura~\ref{sfig:00-alt} destaca as primeiras 20 amostras
            obtidas pela cadeia ao come\'car em $(x=0, y=0)$. Elas
            parecem ser amostras representativas da distribui\'c\~ao;
            no entanto, n\~ao s\~ao suficientes para aproximar a distribui\'c\~ao
            por conta pr\'opria. Isso significaria que uma m\'edia amostral computada
            com apenas 20 amostras teria alta vari\^ancia, i.e.\ seu
            valor dependeria fortemente dos valores das 20 amostras
            usadas para computar a m\'edia.
            \item A Figura~\ref{sfig:77-alt} destaca as primeiras 20 amostras
            obtidas pela cadeia ao come\'car em $(x=7, y=7)$. Pode-se
            claramente ver a cauda de ``burn-in'' que se aproxima lentamente da
            regi\~ao onde se encontra a maior parte da massa de probabilidade.
        \end{itemize}
        
        
    \end{solution}
    
    \item Na pr\'atica, n\~ao sabemos onde a distribui\'c\~ao que desejamos
    amostrar possui alta densidade, por isso costumamos inicializar a
    cadeia de Markov de forma um tanto arbitr\'aria, ou na amostra de m\'aximo a-posteriori (MAP), se
    dispon\'ivel. As amostras obtidas no in\'icio da cadeia s\~ao
    tipicamente descartadas, pois n\~ao s\~ao consideradas representativas
    da distribui\'c\~ao alvo. Este per\'iodo inicial entre a
    inicializa\'c\~ao e o in\'icio da coleta de amostras \'e chamado de
    ``warm-up'' (aquecimento), ou tamb\'em de ``burn-in''.
    
    Estenda sua fun\'c\~ao \lstinline|mh| para incluir um argumento de aquecimento
    adicional $W$, que especifica o n\'umero de passos de MCMC dados antes de
    come\'car a coletar amostras. Sua fun\'c\~ao deve ainda retornar uma lista
    de $S$ amostras como em \ref{qpt:mh-alt}.
    
    \begin{solution}
        
        Podemos estender a fun\'c\~ao \lstinline|mh| com um argumento de aquecimento, por exemplo, iterando por \lstinline|num_samples + warmup|
        passos, e come\'car a registrar amostras apenas ap\'os o per\'iodo de aquecimento:
        \begin{lstlisting}
            def mh(p_star, param_init, num_samples=5000, vari=1.0, warmup=0):
            x = []
            x_current = param_init
            for n in range(num_samples+warmup):
            ...  # corpo igual ao anterior
            
            if n >= warmup: x.append(x_next)
            x_current = x_next
            
            return x
        \end{lstlisting}
    \end{solution}
    
\end{exenumerate}


\ex{Regress\~ao de Poisson Bayesiana}\label{q:poisson-reg-alt}
Considere um modelo de regress\~ao de Poisson
Bayesiano, onde as sa\'idas $y_n$ s\~ao geradas de uma distribui\'c\~ao de Poisson
de taxa $\exp(\alpha x_n + \beta)$, onde os $x_n$ s\~ao as entradas (covari\'aveis), e $\alpha$ e $\beta$ os par\^ametros do modelo de regress\~ao para os quais assumimos uma priori Gaussiana ampla:
\begin{align}
    \alpha &\sim \normal(\alpha; 0, 100) \\
    \beta &\sim \normal(\beta; 0, 100)\\
    y_n &\sim \mathrm{Poisson}(y_n; \exp(\alpha x_n + \beta)) \quad \text{para } n = 1, \dots, N 
\end{align}
$\mathrm{Poisson}(y ;\lambda)$ denota a fun\'c\~ao de massa de probabilidade
de uma vari\'avel aleat\'oria de Poisson com taxa $\lambda$,
\begin{equation}
    \mathrm{Poisson}(y; \lambda) = \frac{\lambda^y}{y!}\exp(-\lambda), \quad \quad y \in \{0, 1, 2, \ldots\}, \quad \lambda > 0
\end{equation}
Considere $\data = \{(x_n, y_n)\}_{n=1}^N$ onde $N=5$ e:
\begin{align}
    (x_1, \ldots, x_5) &= (-0.50519053, -0.17185719,  0.16147614,  0.49480947,  0.81509851)\\
    (y_1, \ldots, y_5) &= (1, 0, 2, 1, 2)
\end{align}

Estamos interessados em computar a densidade posterior dos par\^ametros
$(\alpha, \beta)$ dados os dados $\data$ acima.

\begin{exenumerate}
    
    \item Derive uma express\~ao para a densidade posterior n\~ao normalizada de
    $\alpha$ e $\beta$ dados $\data$, i.e. uma fun\'c\~ao $p^*$ dos
    par\^ametros $\alpha$ e $\beta$ que \'e proporcional \`a
    densidade posterior $p(\alpha, \beta\mid \data)$, e que pode assim
    ser usada como densidade alvo no algoritmo Metropolis Hastings.
    
    \begin{solution}
        Pela regra do produto, a distribui\'c\~ao conjunta descrita pelo
        modelo, com os dados $\data$ inseridos, \'e proporcional \`a posterior
        e, portanto, pode ser tomada como $p^*$:
        \begin{align}
            p^*(\alpha, \beta) &= p(\alpha, \beta, \{(x_n,y_n)\}_{n=1}^N ) \\ 
            &= \normal(\alpha; 0, 100)\normal(\beta; 0, 100)\prod_{n=1}^{N}\mathrm{Poisson}(y_n \mid \text{exp}(\alpha x_n + \beta))
        \end{align}
    \end{solution}
    
    \item Implemente a densidade posterior n\~ao normalizada $p^*$ derivada. Se
    seu ambiente de codifica\'c\~ao fornecer uma implementa\'c\~ao da fmp de Poisson
    acima, voc\^e pode us\'a-la diretamente em vez de implementar a
    fmp por conta pr\'opria.
    
    Use o algoritmo Metropolis Hastings da Quest\~ao
    \ref{q:basic-mcmc-alt}\ref{qpt:mh_test-alt} para sortear $5.000$
    amostras da densidade posterior $p(\alpha, \beta\mid \data)$. Defina os
    hiperpar\^ametros do algoritmo Metropolis-Hastings como:
    \begin{itemize}
        \item \lstinline{param_init} $=(\alpha_{\mathrm{init}},\beta_{\mathrm{init}}) = (0,0)$,
        \item \lstinline{vari} $= 1$, e 
        \item n\'umero de passos de aquecimento $W = 1000$.
    \end{itemize}
    
    Plote as amostras colhidas com o eixo x sendo $\alpha$ e o eixo y sendo $\beta$, e relate a m\'edia posterior de $\alpha$ e $\beta$, bem como seu coeficiente de correla\'c\~ao sob a posterior.
    
    \begin{solution}
        Uma implementa\'c\~ao em Python \'e:
        
        \begin{lstlisting}
            import numpy as np
            from scipy.stats import multivariate_normal, norm, poisson, uniform
            
            xx1 = np.array([-0.5051905265552105, -0.17185719322187715, 0.16147614011145617, 0.49480947344478954, 0.8150985069051909])
            yy1 = np.array([1, 0, 2, 1, 2])
            N1 = len(xx1)
            
            def poisson_regression(params):
            a = params[0]
            b = params[1]
            # media zero, desvio padrao 10 == variancia 100
            p = norm.pdf(a, loc=0, scale=10) * norm.pdf(b, loc=0, scale=10) 
            for n in range(N1):
            p = p * poisson.pmf(yy1[n], np.exp(a * xx1[n] + b)) 
            
            return p
            
            # amostrar
            S = 5000
            samples = np.array(mh(poisson_regression, np.array([0, 0]), num_samples=S, vari=1.0, warmup=1000))
        \end{lstlisting}
        
        Um gr\'afico de dispers\~ao mostrando $5.000$ amostras da posterior \'e exibido na
        Figura~\ref{fig:easy-alt}. A m\'edia posterior de $\alpha$ \'e 0.84, a
        m\'edia posterior de $\beta$ \'e -0.2, e o coeficiente de correla\'c\~ao
        posterior \'e -0.63. Note que os valores num\'ericos s\~ao espec\'ificos para cada amostra.
        \begin{figure}
            \centering
            \includegraphics[width=0.5\textwidth]{poisson-posterior}
            \caption{Amostras posteriores para o problema de regress\~ao de Poisson;
                $\thetab_{\mathrm{init}} = (0,0)$.}
            \label{fig:easy-alt}
        \end{figure}
        
    \end{solution}
    
    
\end{exenumerate}


\ex{Mistura e converg\^encia de MCMC Metropolis-Hasting}\label{q:mh_mixing-alt} Sob condi\'c\~oes fracas, um
algoritmo MCMC \'e um algoritmo de infer\^encia assintoticamente exato,
o que significa que, se for executado para sempre, gerar\'a amostras que
correspondem \`a distribui\'c\~ao de probabilidade desejada. Neste caso, diz-se
que a cadeia converge.

Na pr\'atica, queremos executar o algoritmo por tempo suficiente para sermos
capazes de aproximar a posterior adequadamente. O que significa tempo suficiente para
que a cadeia convirja varia drasticamente dependendo do algoritmo, dos
hiperpar\^ametros (e.g.\ a vari\^ancia \lstinline{vari}) e da distribui\'c\~ao
posterior alvo. \'E imposs\'ivel determinar exatamente se a cadeia foi
executada por tempo suficiente, mas existem v\'arios diagn\'osticos que podem
nos ajudar a determinar se podemos ``confiar'' na aproxima\'c\~ao baseada em
amostras para a posterior.

Uma maneira muito r\'apida e comum de avaliar a converg\^encia da cadeia de Markov
\'e inspecionar visualmente os \emph{gr\'aficos de tra\'co} (trace plots) para cada
par\^ametro. Um gr\'afico de tra\'co mostra como as amostras colhidas evoluem ao
longo do tempo, ou seja, s\~ao uma s\'erie temporal das amostras geradas pela
cadeia de Markov. A Figura~\ref{fig:traceplots-alt} mostra exemplos de gr\'aficos de tra\'co
obtidos pela execu\'c\~ao do algoritmo Metropolis Hastings para diferentes valores
dos hiperpar\^ametros \lstinline{vari} e \lstinline{param_init}. Idealmente, a
s\'erie temporal cobre todo o dom\'inio da distribui\'c\~ao alvo e \'e dif\'icil
``ver'' qualquer estrutura nela, de modo que prever valores de amostras futuras a
partir da atual seja dif\'icil. Se for assim, as amostras s\~ao provavelmente independentes
umas das outras e diz-se que a cadeia est\'a bem ``misturada'' (mixed).



\begin{exenumerate}
    \item \label{qpt:mh_trace-alt} Considere os gr\'aficos de tra\'co na Figura~\ref{fig:traceplots-alt}: A
    vari\^ancia \lstinline{vari} usada na Figura
    \ref{sfig:traceplot_many_small-alt} \'e maior ou menor que o valor de
    \lstinline{vari} usado na Figura \ref{sfig:traceplot_many_good-alt}? A
    vari\^ancia \lstinline{vari} usada na Figura \ref{sfig:traceplot_many_big-alt}
    \'e maior ou menor que o valor usado na Figura
    \ref{sfig:traceplot_many_good-alt}?
    
    Em ambos os casos, explique o comportamento dos gr\'aficos de tra\'co em termos
    do funcionamento do algoritmo Metropolis Hastings e do efeito da
    vari\^ancia \lstinline{vari}.
    
    \begin{figure}[p]
        \setcounter{subfigure}{0}
        \centering
        \subfloat[vari\^ancia \lstinline{vari}: $1$]{
            \includegraphics[width=\textwidth]{Q3d_traceplot_stepsize1}
            \label{sfig:traceplot_many_good-alt}}\\
        \subfloat[Valor alternativo de \lstinline{vari}]{
            \includegraphics[width=\textwidth]{Q3d_traceplot_stepsize001}
            \label{sfig:traceplot_many_small-alt}}\\
        \subfloat[Valor alternativo de \lstinline{vari}]{
            \includegraphics[width=\textwidth]{Q3d_traceplot_stepsize50}
            \label{sfig:traceplot_many_big-alt}}
        \caption{Para a Quest\~ao \ref{q:mh_mixing-alt}\ref{qpt:mh_trace-alt}: Gr\'aficos de tra\'co do par\^ametro $\beta$ da Quest\~ao \ref{q:poisson-reg-alt} sorteados usando Metropolis-Hastings com diferentes vari\^ancias da distribui\'c\~ao de proposta.}
        \label{fig:traceplots-alt}
    \end{figure}
    %
    
    \begin{solution}
        
        M\'etodos MCMC s\~ao sens\'iveis a diferentes hiperpar\^ametros, e
        geralmente precisamos diagnosticar cuidadosamente os resultados da infer\^encia para garantir
        que nosso algoritmo aproxime adequadamente a distribui\'c\~ao posterior alvo.
        
        
        \begin{enumerate}[label=(\roman*)]
            
            \item A Figura~\ref{sfig:traceplot_many_small-alt} utiliza uma vari\^ancia \textit{pequena}
            (\lstinline{vari} foi definida como $0,001$). Os gr\'aficos de tra\'co mostram
            que as amostras para $\beta$ s\~ao altamente correlacionadas e evoluem
            muito lentamente ao longo do tempo. Isso ocorre porque a aleatoriedade
            introduzida \'e bastante pequena em compara\'c\~ao com a escala da
            posterior, logo a amostra proposta em cada itera\'c\~ao de MCMC estar\'a
            muito pr\'oxima da amostra atual e, portanto, ser\'a provavelmente aceita.
            
            Explica\'c\~ao mais matem\'atica: para uma distribui\'c\~ao de proposta sim\'etrica, a
            raz\~ao de aceita\'c\~ao $a$ torna-se:
            \begin{equation}
                a =  \frac{p^*(\thetab^*)}{p^*(\thetab)},
            \end{equation}
            onde $\thetab$ \'e a amostra atual e $\thetab^*$ \'e a amostra
            proposta. Para vari\^ancias que s\~ao pequenas em rela\'c\~ao \`a escala (ao quadrado)
            da posterior, $a$ est\'a pr\'oximo de um e a amostra proposta $\thetab^*$
            \'e provavelmente aceita. Isso d\'a origem \`as s\'eries temporais de mudan\'ca lenta
            mostradas na Figura~\ref{sfig:traceplot_many_small-alt}.
            
            \item Na Figura~\ref{sfig:traceplot_many_big-alt}, a vari\^ancia \'e
            maior que a de refer\^encia (\lstinline{vari} foi definida como $50$). Os gr\'aficos
            de tra\'co sugerem que muitas itera\'c\~oes do algoritmo resultam na
            rejei\'c\~ao da amostra proposta e, assim, acabamos copiando a
            mesma amostra repetidamente. Isso ocorre porque, se as perturba\'c\~oes aleat\'orias
            forem grandes em rela\'c\~ao \`a escala da
            posterior, $p^*(\thetab^*)$ pode ser muito diferente de
            $p^*(\thetab)$ e $a$ pode ser muito pequeno.
            
        \end{enumerate}
        
    \end{solution}
    
    \item No Metropolis-Hastings, e no MCMC em geral, qualquer amostra depende
    da amostra gerada anteriormente e, portanto, o algoritmo
    gera amostras que s\~ao geralmente estatisticamente dependentes. O
    \textit{tamanho efetivo da amostra} de uma sequ\^encia de amostras dependentes \'e
    o n\'umero de amostras independentes que s\~ao, em algum sentido,
    equivalentes ao nosso n\'umero de amostras dependentes. Uma defini\'c\~ao do
    tamanho efetivo da amostra (ESS - Effective Sample Size) \'e:
    \begin{equation}
        \text{ESS} = \frac{S}{1+ 2\sum_{k=1}^{\infty} \rho(k)}
    \end{equation}
    onde $S$ \'e o n\'umero de amostras dependentes colhidas e $\rho(k)$ o
    coeficiente de correla\'c\~ao entre duas amostras na cadeia de Markov que
    est\~ao a $k$ pontos de tempo de dist\^ancia. Podemos ver que, se as amostras forem
    fortemente correlacionadas, $\sum_{k=1}^{\infty} \rho(k)$ \'e grande e o
    tamanho efetivo da amostra \'e pequeno. Por outro lado, se $\rho(k)=0$
    para todo $k$, o tamanho efetivo da amostra \'e $S$.
    
    $\text{ESS}$, conforme definido acima, \'e o n\'umero de amostras independentes
    necess\'arias para obter uma m\'edia amostral que possua a mesma
    vari\^ancia que a m\'edia amostral computada a partir de amostras correlacionadas.
    
    Para ilustrar como a correla\'c\~ao entre amostras est\'a relacionada a
    uma redu\'c\~ao do tamanho da amostra, considere dois pares de amostras $(\theta_1,
    \theta_2)$ e $(\omega_1, \omega_2)$. Todas as vari\'aveis possuem vari\^ancia
    $\sigma^2$ e a mesma m\'edia $\mu$, mas $\omega_1$ e $\omega_2$ s\~ao
    n\~ao correlacionadas, enquanto a matriz de covari\^ancia para $\theta_1, \theta_2$ \'e
    $\C$:
    \begin{equation}
        \C = \sigma^2 \begin{pmatrix}
            1 & \rho \\
            \rho & 1
        \end{pmatrix},
    \end{equation}
    com $\rho>0$. A vari\^ancia da m\'edia $\bar{\omega} = 0.5
    (\omega_1+\omega_2)$ \'e:
    \begin{equation}
        \var \left( \bar{\omega} \right) = \frac{\sigma^2}{2},
    \end{equation}
    onde o $2$ no denominador \'e o tamanho da amostra.
    
    Derive uma equa\'c\~ao para a vari\^ancia de $\bar{\theta} =
    0.5(\theta_1+\theta_2)$ e compute a redu\'c\~ao $\alpha$ do
    tamanho da amostra ao trabalhar com as correlacionadas $(\theta_1,\theta_2)$. Em
    outras palavras, derive uma equa\'c\~ao de $\alpha$ em:
    \begin{equation}
        \var \left(\bar{\theta}\right) = \frac{\sigma^2}{2/\alpha}.
    \end{equation}
    Qual \'e o tamanho efetivo da amostra $2/\alpha$ quando $\rho \to 1$?
    
    \begin{solution}
        
        Note que $\E(\bar{\theta}) = \mu$. Da defini\'c\~ao de
        vari\^ancia, temos ent\~ao:
        \begin{align}
            \var(\bar{\theta}) & = \E\left( (\bar{\theta} - \mu)^2\right)\\
            & = \E\left( \left(\frac{1}{2}(\theta_1+\theta_2) - \mu\right)^2\right) \\
            & =  \E\left( \left(\frac{1}{2}(\theta_1-\mu+\theta_2-\mu)\right)^2\right) \\
            & =  \frac{1}{4} \E\left( (\theta_1-\mu)^2+(\theta_2-\mu)^2 + 2 (\theta_1-\mu) (\theta_2-\mu)\right) \\
            & =  \frac{1}{4} ( \sigma^2 + \sigma^2 + 2 \sigma^2 \rho) \\
            & =  \frac{1}{4} ( 2\sigma^2 + 2 \sigma^2 \rho) \\
            & =  \frac{\sigma^2}{2} ( 1 + \rho) \\
            & =  \frac{\sigma^2}{2/(1+\rho)}
        \end{align}
        Logo: $\alpha = (1+\rho)$, e para $\rho \to 1$, $2/\alpha \to 1$.
        
        Devido \`a forte correla\'c\~ao, temos efetivamente apenas uma amostra e n\~ao duas se $\rho \to 1$.
        
    \end{solution}
    
\end{exenumerate}